\documentclass[11pt]{article}
\usepackage[margin=1in]{geometry}
\usepackage{amsmath,amssymb,amsthm}
\usepackage{graphicx}
\usepackage{hyperref}

\newtheorem{theorem}{Theorem}
\newtheorem{corollary}{Corollary}
\newtheorem{definition}{Definition}

\newcommand{\bbT}{\mathbb T}
\newcommand{\Ree}{\operatorname{Re}}
\newcommand{\dist}{\operatorname{dist}}

\title{A Two-Sided Strong-Exposure Obstruction to Numerical Index One}
\author{agent\_lane\_03}
\date{June 23, 2026}

\begin{document}
\maketitle

\section*{Source Problem}

The source paper is Pal--Roy--Sain, \emph{A characterization of Banach spaces with numerical index one}, arXiv:2604.03690 \cite{PRS}.  On page 7, after recalling necessary conditions of Lopez--Martin--Paya \cite{LMP}, the authors write that it is open whether the real reflexive finite-dimensionality theorem remains true for complex Banach spaces:

\begin{center}
\includegraphics[width=\textwidth]{figures/open_problem_crop.png}
\end{center}

In other words, the open problem is whether every complex reflexive Banach space with numerical index one must be finite-dimensional.

\section*{Proof Intuition}

The numerical index condition $n(X)=1$ says that every operator has numerical radius equal to its operator norm.  A two-sided strongly exposed support pair gives enough local rigidity to defeat this for a rank-one operator.  If $x^*$ strongly exposes $x$, then any numerical-radius test vector that nearly norms the rank-one functional factor $x^*$ must be close to the circle orbit of $x$.  If, at the same time, $x$ strongly exposes $x^*$ in the dual ball, then the corresponding support functionals must be close to the circle orbit of $x^*$.  Choosing the rank-one operator $Tz=x^*(z)y$ with $y\in \ker x^*$ forces those same support functionals to be almost zero on $y$, contradicting numerical radius one.

Thus a complex reflexive counterexample, if it exists, cannot have even one such two-sided smooth/rotund support point.  The remaining difficulty in the full open problem is that reflexivity supplies denting phenomena, but not automatically this matched two-sided strong exposure for the original norm.

\section*{Definitions}

Let $X$ be a complex Banach space.  The numerical radius of $T\in L(X)$ is
\[
  w(T)=\sup\{|x^*(Tx)|:\ x\in S_X,\ x^*\in S_{X^*},\ x^*(x)=1\},
\]
and the numerical index is
\[
  n(X)=\inf\{w(T):\ T\in L(X),\ \|T\|=1\}.
\]

\begin{definition}
Let $x\in S_X$ and $x^*\in S_{X^*}$ satisfy $x^*(x)=1$.
We say that $(x,x^*)$ is a two-sided strongly exposed support pair if:
\begin{enumerate}
\item whenever $z_\alpha\in B_X$ and $\Ree x^*(z_\alpha)\to 1$, then $\|z_\alpha-x\|\to 0$;
\item whenever $f_\alpha\in B_{X^*}$ and $\Ree f_\alpha(x)\to 1$, then $\|f_\alpha-x^*\|\to 0$.
\end{enumerate}
\end{definition}

The first condition says that $x^*$ strongly exposes $B_X$ at $x$.  The second says that $x$, viewed as an element of $X^{**}$, strongly exposes $B_{X^*}$ at $x^*$.

\section*{Main Partial Result}

\begin{theorem}
Let $X$ be a complex Banach space with $\dim X\ge 2$.  If $X$ has a two-sided strongly exposed support pair, then $n(X)<1$.
\end{theorem}

\begin{proof}
Let $(x,x^*)$ be such a pair.  Since $\dim X\ge 2$, choose $y\in S_X\cap \ker x^*$.  Define the rank-one operator
\[
  Tz=x^*(z)y\qquad (z\in X).
\]
Then $\|T\|=1$.  We claim that $w(T)<1$.

Suppose, toward a contradiction, that $w(T)=1$.  By the definition of the supremum, for each $n$ there are $u_n\in S_X$ and $u_n^*\in S_{X^*}$ such that $u_n^*(u_n)=1$ and
\[
  |u_n^*(Tu_n)|=|x^*(u_n)|\,|u_n^*(y)|\longrightarrow 1.
\]
Hence both $|x^*(u_n)|\to 1$ and $|u_n^*(y)|\to 1$.

Put $\lambda_n=\overline{x^*(u_n)}/|x^*(u_n)|$ after discarding finitely many terms for which $x^*(u_n)=0$, and set
\[
  v_n=\lambda_n u_n,\qquad v_n^*=\overline{\lambda_n}u_n^*.
\]
Then $v_n^*(v_n)=1$, $|v_n^*(y)|=|u_n^*(y)|\to 1$, and
\[
  x^*(v_n)=|x^*(u_n)|\to 1.
\]
Since $x^*$ strongly exposes $B_X$ at $x$, we have $\|v_n-x\|\to 0$.
Consequently,
\[
  |v_n^*(x)-1|=|v_n^*(x)-v_n^*(v_n)|\le \|x-v_n\|\to 0.
\]
Thus $\Ree v_n^*(x)\to 1$.  Since $x$ strongly exposes $B_{X^*}$ at $x^*$, it follows that $\|v_n^*-x^*\|\to 0$.  But $x^*(y)=0$, so
\[
  |v_n^*(y)|\le \|v_n^*-x^*\|\|y\|+|x^*(y)|\to 0,
\]
contradicting $|v_n^*(y)|\to 1$.  Therefore $w(T)<1$, and since $\|T\|=1$, we get $n(X)\le w(T)<1$.
\end{proof}

\begin{corollary}
Let $X$ be a complex Banach space with $n(X)=1$ and $\dim X\ge 2$.  Then $X$ has no two-sided strongly exposed support pair.
\end{corollary}

\begin{corollary}
If a complex Banach space $X$ of dimension at least two has a point $x\in S_X$ which is locally uniformly rotund and Frechet smooth, and if its unique support functional $x^*$ satisfies the corresponding dual strong-exposure condition above, then $n(X)<1$.  In particular, a uniformly convex and uniformly smooth complex Banach space of dimension at least two has numerical index strictly smaller than one.
\end{corollary}

\begin{proof}
Local uniform rotundity at $x$ implies that the unique support functional $x^*$ strongly exposes $B_X$ at $x$.  Frechet smoothness at $x$ is the standard duality condition giving strong exposure of $B_{X^*}$ at $x^*$.  Apply the theorem.  The uniformly convex and uniformly smooth case satisfies these hypotheses at every unit vector.
\end{proof}

\section*{Scope and Limitations}

This is a partial obstruction, not a solution to the full open problem.  It proves that any infinite-dimensional complex reflexive Banach space with numerical index one would have to avoid matched two-sided strong exposure in its original norm.  Reflexive spaces have the Radon--Nikodym property and denting structure, but denting alone is weaker than the two-sided strong exposure used in the rank-one argument above.

\section*{Novelty and Search Notes}

The lightweight run indexes had no entry for arXiv:2604.03690 or this two-sided strong-exposure obstruction.  Web searches for the exact complex-reflexive numerical-index-one formulation mainly returned the source paper itself.  Nearby run memory contains numerical-radius and Daugavet/lushness packets, but none addresses this source problem or this rank-one obstruction.

\section*{Verification Notes}

The proof is a direct contradiction argument with the rank-one operator $Tz=x^*(z)y$.  The only external facts used are the source formulation of the open problem and standard duality between local uniform rotundity/Frechet smoothness and strong exposure.  No numerical computation is involved.  The human reviewer should mainly check the phase normalization in the complex numerical-radius argument and the exact hypotheses in the corollaries.

\begin{thebibliography}{9}

\bibitem{PRS}
S. Pal, S. Roy, and D. Sain,
\emph{A characterization of Banach spaces with numerical index one},
arXiv:2604.03690, 2026.

\bibitem{LMP}
G. Lopez, M. Martin, and R. Paya,
\emph{Real Banach spaces with numerical index 1},
Bulletin of the London Mathematical Society 31 (1999), 207--212.

\end{thebibliography}

\end{document}

